\Qhead{2.9}{Стационарные волны и солитоны. Взаимодействия плоских волн в диспергирующих средах}{Уравнение Кортевега-де-Вриза и синус - Гордона. Стационарные волны. Понятие о солитонах. Взаимодействия плоских волн в диспергирующих средах. Генерация второй гармоники. Параметрическое усиление и генерация.}
\Subfile{2.9.1}{sections/02_waves/q09/s01}
\Subfile{2.9.2}{sections/02_waves/q09/s02}

\recap{Дисперсия может компенсировать нелинейность и дать солитон. Резонансное смешение требует сохранения частоты и согласования фаз; Мэнли–Роу связывает обмен энергией трёх волн.}
\board{Стационарное уравнение КдВ и sech$^2$; кинк; три амплитуды; $\operatorname{sinc}^2$ при фазовой расстройке.}
\oral{Чем баланс в солитоне отличается от баланса в ударном фронте Бюргерса?}
\bait{Можно оставить длину когерентности: после накопления фазовой ошибки $\pi$ следующая часть среды начинает частично отменять уже созданную гармонику.}
\sources{Цысарь~\cite{tsysar2}, лекции~3--8, PDF-с.~15--36; Белокопытов~\cite{belokopytov2}, параметрическое взаимодействие. Синус–Гордон — стандартная отдельная модель; формула кинка проверена подстановкой.}
